%قبل از شروع مسئله ترجیحاً چیزی ننویسید. اما اگر خواستید بنویسید، در قالب لیستی (فقرات و شمارش) بنویسید.
%برای اضافه کردن نمره، بعدی از شروع‌مسئله، [x] را اضافه کنید و به جای x  نمره مورد نظر را قرار دهید. برای زیرمسئله‌ها هم به همین صورت است.
%\شروع‌مسئله[10]

\شروع‌مسئله[]
داده ساختار اصلی یک پشته برای عملیات‌های push و pop است. برای یافتن min نیز از یک پشته‌ی کمکی هم ارتفاع با پشته‌ی اصلی استفاده می‌کنیم که در ارتفاع $h$م آن، کمینه‌ی $h$ عنصر اول پشته‌ی اصلی را نگه می‌داریم. برای مدیریت این استک نیز به صورت زیر عمل می‌کنیم:

هنگام pop از استک اصلی، از این استک نیز pop می‌کنیم.
هنگام push کردن به استک اصلی، عضو جدید را با اخرین عضو پشته  کمکی مقایسه می‌کنیم. اگر کوچکتر بود، عضو جدید وگرنه آخرین عضو پشته‌ی کمکی را به پشته‌ی کمکی push می‌کنیم. از آنجایی که عضو انتهایی کمینه‌ی سایر اعضا تا به اینجا بوده‌است، با روشی که ارائه کردیم، عضو جدید پشته‌ی کمکی، کمینه‌ی اعضای پشته‌ی اصلی می‌باشد.

برای
\متن‌لاتین{deleteMin}،
نمی‌توان این کار را انجام داد. برای اثبات فرض کنید که بشود deleteMin را در مرتبه‌ی $O(1)$ انجام داد.
 در این صورت تمام اعضای یک ارایه را در $O(n)$ درون این داده ساختار می‌ریزیم سپس یکی یکی کمینه را یافته در یک آرایه‌ی جدید ریخته و کمینه را از داده‌ساختار فعلی حذف کند تا داده ساختار خالی شود و هر کدام از این عملیات‌ها در $O(1)$ و در کل در $O(n)$ انجام می‌شود. در این صورت
توانسته‌ایم که یک آرایه‌ی دلخواه را در زمان $O(n)$ مرتب کنیم که طبق قضایا می‌دانیم که مرتب‌سازی آرایه حداقل از مرتبه‌ی زمانی
$\Omega(n\log(n))$
می‌باشد و این تناقض می‌باشد.

%اگر مسئله چند قسمت داشت به این صورت عمل کنید:
%\شروع{زیرمسائل}
%\زیرمسئله[0] % مسئلۀ ۱آ
%\زیرمسئله % مسئلۀ ۱ب
%\زیرمسئله % مسئلۀ ۱پ
%\زیرمسئله % مسئلۀ ۱ت
%\پایان{زیرمسائل}

\solution{
	پاسخ نمونه.
}




%%%%%%%%%%%%%%%%%%%%%%%%%%%%%%%%%%%%%%%%%%%%%%%%%%%%%
% برخی دستورات مفید برای تمرین‌ها. %
%%%%%%%%%%%%%%%%%%%%%%%%%%%%%%%%%%%%%%%%%%%%%%%%%%%%%

% فرمول‌های نمونه:
%$f(x) = \Theta(x)$
%$T(x, y) \leq \log(x) + 2^y + \binom{2n}{n}$
% {\tt code\_function}


% لیست موارد بدون شماره‌گذاری:
%\شروع{فقرات}
%    \فقره مورد اول 
%        \شروع{فقرات}
%            \فقره مورد اول داخلی 
%                \شروع{فقرات}
%                    \فقره اولین مورد خیلی داخلی 
%                    \فقره دومین مورد خیلی داخلی 
%                \پایان{فقرات}
%            \فقره مورد دوم داخلی 
%        \پایان{فقرات}
%    \فقره مورد دوم 
%\پایان{فقرات}

% لیست موارد شماره‌گذاری شده:
%\شروع{شمارش}
%    \فقره یک مورد 
%    \فقره یک مورد دیگر 
%    \فقره باز هم یک مورد دیگر
%\پایان{شمارش}

% برای نوشتن یک فرمول چند خطی که یک شماره بخوره. برای حالت‌های دیگر به سایت زیر مراجعه کنید:
% https://www.overleaf.com/learn/latex/Aligning_equations_with_amsmath
%\begin{align} 
%    \begin{split}
%        (x+y)^3 &= (x+y)^2(x+y) \\
%                &= (x^2+2xy+y^2)(x+y) \\
%                &= (x^3+2x^2y+xy^2) + (x^2y+2xy^2+y^3) \\
%                &= x^3+3x^2y+3xy^2+y^3
%    \end{split}                                 
%\end{align}

% نوشتن ماتریس و آرایه دوبعدی
%\begin{align}
%    A = 
%    \begin{bmatrix}
%        A_{11} & A_{21} \\
%        A_{21} & A_{22}
%    \end{bmatrix}
%\end{align}

% اضافه کردن تصویر با شرح. برای اضافه کردن تصویر بدون شرح، فقط دستور درج‌تصویر را استفاده کنید.
%\شروع{شکل}[h!]
%\تنظیم‌ازوسط
%  \درج‌تصویر[پهنا=0.5\پهنای‌متن]{img.jpg}
%\شرح{شرح تصویر}
%\برچسب{تصویر:۱}
%\پایان{شکل}
%\پایان{زیرمسائل}

%کد زیر را ببینید.
%%\lstinputlisting{source_filename.py}
%\begin{latin}
%\begin{lstlisting}
%# Import the modules
%import sys
%import random
%
%ans = True
%
%while ans:
%question = raw_input("Ask the magic 8 ball a question: (press enter to quit) ")
%
%answers = random.randint(1,8)
%
%if question == "":
%sys.exit()
%
%elif answers == 1:
%print "It is certain"
%
%elif answers == 2:
%print "Outlook good"
%
%elif answers == 3:
%print "You may rely on it"
%
%elif answers == 4:
%print "Ask again later"
%
%elif answers == 5:
%print "Concentrate and ask again"
%
%elif answers == 6:
%print "Reply hazy, try again"
%
%elif answers == 7:
%print "My reply is no"
%
%elif answers == 8:
%print "My sources say no"
%\end{lstlisting}
%\end{latin}